% !TEX TS-program = lualatex
% !TEX encoding = UTF-8

\documentclass[VesperaleOP.tex]{subfiles}

\ifcsname preamble@file\endcsname
  \setcounter{page}{\getpagerefnumber{M-vop32_commune_non_virginum}}
\fi

\begin{document}

\festum{}{Commune\\unius non Virginis}{extra tempus paschale}{Commune non Virginis extra t. p.}{1}
\addcontentsline{toc}{section}{Commune non Virginis extra t. p.}

\lineaparva

\titulumrubrica{Ad II Vesperas}
\cantus[Ant. V]{CNVEXVA}{Super Ps.}
\begin{psalmodia}
\psalmusinpsalmodia{109}
\psalmusinpsalmodia{112}
\psalmusinpsalmodia{121}
\psalmusinpsalmodia{126}
\psalmusinpsalmodia{147}
\end{psalmodia}

\rubrica{Pro Martyre tantum:}
\capitulum{Eccli. 51}{Confitébor tibi, Dómine rex, et collaudábo te Deum salvatórem meum.\pscross{}Confitébor nómini tuo, quóniam adiútor et protéctor factus es mihi,\psstar{} et liberásti corpus meum a perditióne.\pscross{}}

\rubrica{Pro nec Virgine nec Martyre:}
\capitulum{Eccli. 51}{Confitébor tibi, Dómine rex, et collaudábo te Deum salvatórem meum.\pscross{}Confitébor nómini tuo, quóniam adiútor et protéctor factus es mihi,\psstar{} et liberásti corpus meum a perditióne.\pscross{}}

\rubrica{In festo II classis:}
\cantus{CNVEXVH1}{Hymnus}

\rubrica{In festis III classis ubi psalmi Dominici dicuntur ad Laudes:}
\cantus{CNVEXVH2}{Hymnus}

\rubrica{In festis III classis ubi psalmi feriales dicuntur ad Laudes:}
\cantus{CNVEXVH3}{Hymnus}

\versiculus{Diffúsa est grátia in lábiis tuis.}{Proptérea benedíxit te Deus in ætérnum.}

\cantus{CNVEXVAM}{Ad Magn.}

\end{document}
